\documentclass[11pt,a4paper]{article}

\usepackage[utf8]{inputenc}
\usepackage[T1]{fontenc}
\usepackage{lmodern}
\usepackage[english]{babel}
\usepackage{amsmath,amssymb}
\usepackage{graphicx}
\usepackage{booktabs}
\usepackage{caption}
\usepackage{siunitx}
\usepackage[hidelinks]{hyperref}
\usepackage[round]{natbib}
\usepackage{geometry}
\geometry{margin=2.6cm}
\usepackage{enumitem}
\usepackage{microtype}

\graphicspath{{../plots/}}
\captionsetup{font=small,labelfont=bf}
\setlength{\parskip}{0.2em}

\title{\vspace{-1.2cm}\bfseries Energy Poverty and the Targeting of Electricity
Subsidies: Evidence from the FOSE in Peru}
\author{Carlo Vilches Cevallos\thanks{Universitat de Barcelona, Spain.
Master's in Economics, Regulation and Competition in Public Services.
Contact: \texttt{cvil\_49@mail.com}. All errors are the author's own.}\\[2pt]
\small Universitat de Barcelona}
\date{June 2026}

\begin{document}
\maketitle

\begin{abstract}
\noindent This paper evaluates whether the 140~kWh monthly consumption
threshold used to determine eligibility for Peru's Electricity Social
Compensation Fund (\textit{Fondo de Compensaci\'on Social El\'ectrica}, FOSE)
adequately identifies energy-poor households. Using microdata from the 2025
Residential Energy Consumption and Use Survey (ERCUE, 14{,}815 households), I
compare the incidence of energy poverty across households near the threshold
using Coarsened Exact Matching (CEM) and inverse probability weighting (IPW),
with two complementary indicators: energy overload (2M) and hidden energy
poverty (HEP). The current threshold exhibits a large inclusion error
(69.4\%): on average there is no significant difference in energy poverty
between comparable eligible and ineligible households (ATT $=-2.37$~pp,
$p=0.330$). However, among non-extreme poor households a sizeable effect
emerges ($-20.29$~pp, $p=0.014$), robust across specifications.
Counterfactual simulations indicate that the FOSE plays a protective role in
rural areas. The findings suggest that the consumption-based design is poorly
targeted but retains positive effects in specific segments, pointing to the
need to revise the eligibility criteria.

\medskip
\noindent\textbf{Keywords:} energy poverty; FOSE; targeting; electricity
subsidy; Peru; Coarsened Exact Matching.\\
\textbf{JEL:} Q41, Q48, H23, I38, L94.
\end{abstract}

\section{Introduction}

Energy poverty---a household's inability to access essential energy services
at an affordable cost---is a policy problem that extends well beyond developed
countries \citep{epov2023}. In Latin America, where electricity coverage
approaches universality but affordability remains precarious, the phenomenon
has received comparatively little empirical attention. Peru illustrates this
tension: according to ERCUE 2025, 97\% of households are connected to the
grid, yet the median energy burden reaches 5.20\% among extreme-poor
households, and the gap between median urban (92~kWh/month) and rural
(16~kWh/month) consumption reflects two radically different energy realities.

Peru's main direct consumer-protection instrument is the FOSE, a cross-subsidy
created in 2001 that grants discounts on the electricity bill to households
consuming below a regulated threshold---currently 140~kWh/month, raised from
100~kWh after the 2022 reform (D.S. 023-2022-PCM). Like many block-tariff
subsidies, the FOSE uses electricity consumption as its sole eligibility
criterion, implicitly assuming that low consumption signals economic
vulnerability. This premise has been questioned across contexts.
\citet{whittington2015} show that consumption-based subsidies are poorly
targeted toward the poorest, and \citet{levinson2022} document that the weak
correlation between electricity consumption and income limits the
redistributive capacity of tariff discounts. In Spain, evaluations of the
analogous \textit{Bono Social El\'ectrico} find modest effects: a reduction of
about two percentage points in the probability of energy poverty
\citep{bagnoli2022}, no significant effect on subjective indicators
\citep{garcia2021}, and a declining effect that is nearly null for the poorest
households \citep{llorca2025}.

Despite the FOSE's relevance for Peruvian social policy, no empirical test of
its eligibility rule or of its effect on household-level energy-poverty
indicators exists, to the best of my knowledge. This paper fills that gap. The
central hypothesis is twofold. \textit{First}, electricity consumption is an
imperfect proxy for vulnerability: in the aggregate, the 140~kWh rule does not
cleanly separate vulnerable from non-vulnerable households because the eligible
group is heavily diluted by non-poor households; yet in the subgroup of
non-extreme poor households the threshold does discriminate. \textit{Second},
even where the rule points in the right direction, the size of the tariff
benefit is too small to produce a clear separation in energy well-being around
the threshold.

The analysis yields three main findings. The FOSE exhibits a 69.4\% inclusion
error and a normative exclusion pattern that disproportionately affects rural
households with shared electricity connections. The 140~kWh threshold produces
no significant separation in the probability of energy overload between
comparable eligible and excluded households in the narrow $\pm20$~kWh window,
though the effect strengthens in wider windows and becomes large and
significant once the sample is restricted to poor households. Finally, hidden
energy poverty (HEP) is essentially zero in the neighbourhood of the
threshold, revealing that energy poverty there manifests exclusively as
financial overload rather than deprivation.

\section{The FOSE and Peru's Energy Context}

Peru has achieved near-universal electrification (97\% of households).
However, grid access guarantees neither adequate consumption nor an affordable
price: the median energy burden is 3.13\% nationally but reaches 5.20\% among
extreme-poor households. The distribution of consumption is bimodal by area
(Figure~\ref{fig:dist}): rural density concentrates below 30~kWh, while the
urban distribution peaks at 60--90~kWh. Overall, 76\% of households consume at
or below the 140~kWh threshold and are therefore eligible for the FOSE.

\begin{figure}[t]
\centering
\includegraphics[width=0.92\textwidth]{fig1_consumption_distribution.png}
\caption{Weighted distribution of monthly electricity consumption and the FOSE
eligibility threshold, ERCUE 2025. Source: ERCUE 2025; weighted by expansion
factor; 5~kWh bins.}
\label{fig:dist}
\end{figure}

The FOSE was created in 2001 (Law 27510) and operates as an intra-sectoral
cross-subsidy: users with consumption above the threshold and free-market
clients finance, through a surcharge, the discounts applied to beneficiaries.
As of December 2024, two of every three residential users were FOSE
beneficiaries, and the annual subsidy reached roughly 540~million soles (an
80\% increase relative to 2019). The 2022 reform (Law 31429) raised the
threshold and introduced additional exclusion criteria: no more than one
property or supply point, not being located in a high- or middle-income block,
and no consumption above 140~kWh during the summer months.

Heterogeneity by typical electricity sector is pronounced
(Table~\ref{tab:sector}). Lima (ST1) records average consumption of 160~kWh
and 42.8\% of households above the threshold, against the rural isolated
systems (SER) with 37.5~kWh and barely 3\%. The discount structure operates in
three blocks: for the first 30~kWh the discount is proportional to the energy
charge; between 31 and 140~kWh it is a fixed amount; above 140~kWh the
household becomes a net contributor to the fund.

\begin{table}[t]
\centering
\small
\caption{Electricity consumption distribution by typical sector, ERCUE 2025.}
\label{tab:sector}
\begin{tabular}{lccccccc}
\toprule
Sector & Mean & Median & $\leq$30 & 31--140 & $>$140 & Poverty \\
 & (kWh) & (kWh) & (\%) & (\%) & (\%) & (\%) \\
\midrule
ST1 (Lima)        & 160.2 & 129.8 & 3.3  & 53.9 & 42.8 & 16.1 \\
ST2 (Urban mid)   & 92.8  & 75.3  & 15.9 & 66.4 & 17.7 & 18.1 \\
ST3 (Urban-rural) & 52.7  & 36.8  & 43.0 & 50.1 & 6.9  & 34.8 \\
ST4 (Rural)       & 32.6  & 17.0  & 67.7 & 29.3 & 3.0  & 32.9 \\
SER (Rural isol.) & 37.5  & 26.9  & 59.5 & 37.5 & 3.0  & 32.3 \\
\bottomrule
\end{tabular}
\\[2pt]
\footnotesize Source: ERCUE 2025.
\end{table}

\section{Data and Empirical Strategy}

\paragraph{Data.} The analysis uses microdata from ERCUE 2025, a probabilistic,
stratified, multi-stage survey representative at the national, departmental and
urban/rural level, with fieldwork conducted between September and November
2024. The sample comprises 15{,}485 households (14{,}815 with valid
consumption data) across 25 departments, and the expansion factor scales the
observations to the universe of 9.3~million households with electricity
access. Two relevant limitations are that ERCUE does not contain the regulated
tariff by electrical system (approximated through the implicit price,
expenditure/consumption) nor the INEI block-stratum classification, and that
the cross-sectional design precludes dynamic analysis.

\paragraph{Energy-poverty indicators.} The household energy burden is
$V_i = G^{\text{elec}}_i / G^{\text{total}}_i$. A household is classified as
energy-overloaded (\textit{2M}) when this ratio exceeds twice the national
median:
\begin{equation}
2M_i = \mathbb{1}\!\left(V_i > 2\widetilde{V}\right).
\end{equation}
Because indicators based solely on relative expenditure can understate
deprivation when households suppress consumption, I add the hidden energy
poverty (HEP) indicator, which flags households whose absolute electricity
expenditure falls below half the median of their reference group
\citep{thomson2016}:
\begin{equation}
\text{HEP}_i = \mathbb{1}\!\left(G^{\text{elec}}_i <
0.5 \times \widetilde{G}^{\text{elec}}_{\text{group}_i}\right),
\end{equation}
where the group is defined by geographic area.

\paragraph{Identification.} The strategy compares households on either side of
the 140~kWh administrative threshold within narrow consumption windows: a main
specification of $\pm20$~kWh ($N\approx1{,}286$) and a robustness window of
$\pm30$~kWh ($N\approx1{,}949$). The parameter of interest is
$\tau = E[Y_i \mid C_i < 140] - E[Y_i \mid C_i \geq 140]$, with
$Y_i \in \{2M_i, \text{HEP}_i\}$; standard errors come from a nonparametric
bootstrap with 2{,}000 replications. Although a regulatory threshold suggests a
regression-discontinuity (RD) design, three features preclude its valid
application: FOSE eligibility does not depend exclusively on consumption, the
data are cross-sectional, and compositional heterogeneity around the threshold
violates the continuity assumption.

\paragraph{Estimation.} Coarsened Exact Matching \citep[CEM,][]{iacus2012} is
the main method. The final specification retains three covariates---grouped
typical sector, area (urban/rural) and household size---generating 17 cells
with 99.8\% common support. The ATT estimator is
$\widehat{\tau}_{\text{ATT}} = \sum_j (n_{1j}/N_1)(\overline{Y_{1j}}-\overline{Y_{0j}})$,
where eligibility is $T_i = \mathbb{1}(C_i<140)$. Inverse probability weighting
(IPW) serves as a sensitivity check. This paper evaluates the discriminating
capacity of the eligibility rule and should not be read as a causal estimate of
the subsidy's impact on welfare.

\section{Results}

\subsection{Targeting errors}

84.1\% of ERCUE 2025 households consume at or below 140~kWh, making them
eligible for the FOSE. This near-universal coverage generates a 69.4\%
inclusion error: seven of every ten consumption-eligible households are not in
monetary poverty. Disaggregation by sector shows the inclusion error is
substantially larger in urban sectors (74.8\% in ST1, 78.6\% in ST2) than in
rural sectors (60--63\%), reflecting greater income heterogeneity in urban
areas \citep{levinson2022}. Exclusion error by the consumption criterion is
5.3\% but rises to 20.3\% in ST1. The main normative exclusion criterion in
rural sectors is not the consumption threshold but the shared-meter condition:
8.5\% of eligible SER households share an electricity supply point, an
exclusion driven by agricultural arrangements rather than ability to pay.

Figure~\ref{fig:cov} shows that the 2022 reform pushed national coverage from
53.8\% to 76.1\% while the inclusion error edged up to 75.7\% under the full
rule. Table~\ref{tab:target} summarises the targeting indicators under
alternative eligibility definitions.

\begin{figure}[t]
\centering
\includegraphics[width=\textwidth]{fig2_coverage_errors.png}
\caption{Evolution of FOSE coverage and targeting errors, 2019--2025. Source:
own elaboration with ERCUE 2019, 2023 and 2025.}
\label{fig:cov}
\end{figure}

\begin{table}[t]
\centering
\small
\caption{Targeting indicators by eligibility definition, ERCUE 2025.}
\label{tab:target}
\begin{tabular}{lccc}
\toprule
Indicator & By consumption & Full rule & Full rule \\
 & & (weighted) & (unweighted) \\
\midrule
Definition       & $C\leq140$ kWh & + meter + stratum & + meter + stratum \\
Coverage         & 84.1\% & 76.1\% & 76.1\% \\
Inclusion error  & 69.4\% & 75.5\% & --- \\
Exclusion error  & 5.3\%  & 10.4\% & 15.2\% \\
\bottomrule
\end{tabular}
\\[2pt]
\footnotesize Source: ERCUE 2025.
\end{table}

\subsection{Discriminating power of the threshold}

The crude difference between eligible and excluded households in the
$\pm20$~kWh window is $-1.03$~pp and not significant. Applying CEM, the ATT is
$-2.37$~pp (SE~$=2.43$; $p=0.330$). The $\pm30$~kWh robustness window yields an
ATT of $-5.72$~pp ($p=0.003$), significant at the 1\% level, and IPW gives a
consistent $-2.09$~pp (SE~$=2.35$; $p=0.376$). Common support is adequate: the
mean propensity score is 0.575 for eligibles and 0.571 for excluded
households. Figure~\ref{fig:est} compares the estimators.

\begin{figure}[t]
\centering
\includegraphics[width=0.92\textwidth]{fig3_att_estimators.png}
\caption{Difference in energy overload (2M) associated with FOSE eligibility:
comparison of estimators. Source: ERCUE 2025.}
\label{fig:est}
\end{figure}

The negative sign of the ATT and the significance in the wide window suggest
the threshold points in the right direction, but the difference is modest and
insignificant in the main window: the FOSE discount is too small to create a
sharp separation. For a household consuming 130~kWh, the median FOSE discount
is S/~10.72, barely 10\% of a S/~105 bill---convergent with international
evidence on tariff subsidies of insufficient scale
\citep{bagnoli2022,garcia2021,llorca2025}. Importantly, in the $\pm20$~kWh
window the HEP rate is exactly 0.0\% for both eligible and excluded households:
near the threshold, energy poverty manifests exclusively as financial overload,
not deprivation.

\subsection{Heterogeneity by poverty status}

The most relevant finding emerges when disaggregating by monetary-poverty
status (Figure~\ref{fig:het}). Among non-extreme poor households in the
$\pm20$~kWh window ($N=121$), the ATT is $-20.29$~pp ($p=0.014$). Among
non-poor households ($N=1{,}132$) the effect is essentially null. For
extreme-poor households ($N=9$) estimation is not feasible due to insufficient
sample. The threshold does possess discriminating power, but only among
households that genuinely face energy vulnerability; because 69.4\% of
eligibles are not poor, the signal dilutes in the full sample. This $-20$~pp
result is stable across all windows from $\pm15$ to $\pm30$~kWh.

\begin{figure}[t]
\centering
\includegraphics[width=0.92\textwidth]{fig4_heterogeneity.png}
\caption{Heterogeneity of the FOSE effect on energy overload by poverty status
(CEM, $\pm20$~kWh). Source: ERCUE 2025.}
\label{fig:het}
\end{figure}

\subsection{Counterfactual: a scenario without the FOSE}

To gauge the magnitude of the FOSE's protective effect, I construct a
counterfactual that removes the discount (Figure~\ref{fig:cf}). The energy
overload rate would rise from 17.0\% to 21.2\% in ST1 ($+4.3$~pp), from 9.7\%
to 18.0\% in ST3 ($+8.3$~pp), and from 12.3\% to 24.3\% in SER ($+12.0$~pp).
The protective effect is largest in rural sectors, where the proportional
discount of the first block (60\%) has greater incidence. Thus, even though the
instrument is poorly targeted in the aggregate, it provides meaningful
protection precisely where consumption is lowest.

\begin{figure}[t]
\centering
\includegraphics[width=0.92\textwidth]{fig5_counterfactual.png}
\caption{Counterfactual simulation: energy overload (2M) with and without the
FOSE, by typical sector. Source: ERCUE 2025; eligible households.}
\label{fig:cf}
\end{figure}

\section{Conclusions and Policy Implications}

This paper has evaluated the targeting and effectiveness of the FOSE as an
energy-poverty mitigation instrument in Peru, using ERCUE 2025 microdata and
quasi-experimental matching methods. Three conclusions follow.

\textit{First}, the FOSE shows structural targeting deficiencies. The 69.4\%
inclusion error is a direct consequence of a quasi-universal design based
solely on electricity consumption---a weak proxy for socioeconomic
vulnerability whose discriminating power varies by sector (74.8--78.6\% in
urban areas versus 60\% in rural areas). The main exclusion criterion in the
SER is normative---the shared-meter rule---not the consumption threshold.

\textit{Second}, the 140~kWh threshold produces no significant separation in
the probability of energy overload between comparable eligible and excluded
households. The null aggregate result is explained by the composition of the
eligible group: 69.4\% are not poor. Among non-extreme poor households the ATT
is $-20.29$~pp ($p=0.014$). The threshold has discriminating power, but it is
masked by the massive inclusion of non-vulnerable households. The
counterfactual shows that removing the FOSE would raise the overload rate by
4.3~pp (ST1) to 12.0~pp (SER).

\textit{Third}, HEP is exactly zero in the analysis window: near the threshold,
energy poverty is financial overload, not deprivation. The FOSE faces two
distinct problems in distinct segments---deprivation among very-low-consumption
households (where the barrier is access) and overload among
intermediate-consumption households (where the barrier is affordability)---and
a single consumption threshold cannot address both at once.

In short, the FOSE does not fail because the subsidy is insufficient, but
because the criterion identifying the beneficiary is weakly correlated with the
dimension of vulnerability it intends to correct. The results suggest three
reform avenues: (i) a transition toward socioeconomic targeting using Peru's
Household Targeting System (SISFOH), following the 2017 reform of Spain's
\textit{Bono Social}; (ii) an urgent review of the multi-meter exclusion rule;
and (iii) a recalibration of the block structure or its complementation with
direct transfers. The main limitation is the cross-sectional design; the
natural extension is a difference-in-differences design exploiting the three
ERCUE editions (2019, 2023 and 2025).

\begin{thebibliography}{99}
\setlength{\itemsep}{2pt}

\bibitem[Arze del Granado et al.(2012)]{arze2012} Arze del Granado, J., Coady,
D., \& Gillingham, R. (2012). The unequal benefits of fuel subsidies: A review
of evidence for developing countries. \textit{World Development, 40}(11),
2234--2248.

\bibitem[Bagnoli \& Bertom\'eu-S\'anchez(2022)]{bagnoli2022} Bagnoli, L., \&
Bertom\'eu-S\'anchez, S. (2022). How effective has the electricity social rate
been in reducing energy poverty in Spain? \textit{Energy Economics, 106},
105792.

\bibitem[Boardman(1991)]{boardman1991} Boardman, B. (1991). \textit{Fuel
poverty: From cold homes to affordable warmth}. Belhaven Press.

\bibitem[Brown et al.(2018)]{brown2018} Brown, C., Ravallion, M., \& van de
Walle, D. (2018). A poor means test? Econometric targeting in Africa.
\textit{Journal of Development Economics, 134}, 109--124.

\bibitem[Coady et al.(2004)]{coady2004} Coady, D., Grosh, M., \& Hoddinott, J.
(2004). Targeting outcomes redux. \textit{The World Bank Research Observer,
19}(1), 61--85.

\bibitem[EPOV(2023)]{epov2023} EPOV (2023). \textit{Energy poverty in the EU}.
EU Energy Poverty Observatory, European Commission.

\bibitem[Garc\'ia \'Alvarez \& Tol(2021)]{garcia2021} Garc\'ia \'Alvarez, G.,
\& Tol, R. S. J. (2021). The impact of the Bono Social de Electricidad on
energy poverty in Spain. \textit{Energy Economics, 103}, 105554.

\bibitem[Iacus et al.(2012)]{iacus2012} Iacus, S. M., King, G., \& Porro, G.
(2012). Causal inference without balance checking: Coarsened exact matching.
\textit{Political Analysis, 20}(1), 1--24.

\bibitem[Levinson \& Silva(2022)]{levinson2022} Levinson, A., \& Silva, E.
(2022). The electric Gini: Income redistribution through energy prices.
\textit{American Economic Journal: Economic Policy, 14}(2), 341--365.

\bibitem[Llorca \& Rodr\'iguez-\'Alvarez(2025)]{llorca2025} Llorca, M., \&
Rodr\'iguez-\'Alvarez, A. (2025). Is Spain's energy voucher lighting the way
for the poor? A microeconomic evaluation of the Bono Social El\'ectrico.
\textit{Cambridge Working Papers in Economics 2542}, University of Cambridge.

\bibitem[Osinergmin(2025)]{osinergmin2025} Osinergmin (2025). \textit{Reporte
de an\'alisis econ\'omico sectorial: Electricidad}. Gerencia de Pol\'iticas y
An\'alisis Econ\'omico.

\bibitem[Thomson et al.(2016)]{thomson2016} Thomson, H., Snell, C., \& Liddell,
C. (2016). Fuel poverty in the European Union: A concept in need of definition?
\textit{People, Place and Policy, 10}(1), 5--24.

\bibitem[Whittington et al.(2015)]{whittington2015} Whittington, D., Nauges,
C., Fuente, D., \& Wu, X. (2015). A diagnostic tool for estimating the
incidence of subsidies delivered by water utilities in low- and medium-income
countries. \textit{Utilities Policy, 34}, 70--81.

\end{thebibliography}

\end{document}
